% GENERATED FILE. Edit canonical lessons, then run npm run generate:books.
% canonical-source-hash: fnv1a64:c798578d6c9ac8e6
% canonical-lessons: FR-C513-orchestre, FR-C513-chef-d-orchestre, FR-C513-partition, FR-C513-repetition, FR-C513-improviser, FR-R513-revision-506-513

\chapter{L'orchestre}
\label{ch:fr-l-orchestre}
\hlchaptermodality{513}

\begin{chapteropening}
\textbf{By the end of this chapter:} \emph{I can talk about playing in an orchestra.}

Complete the last lesson of chapter 513: I can talk about playing in an orchestra.
\end{chapteropening}

\section[l'orchestre]{l'orchestre — orchestra; stalls (in a theatre)}

\label{lesson:FR-C513-orchestre}



\begin{quote}
Before the new one: say the French for chorus, then the French for melody.
\end{quote}

\subsection*{You'll want to know: l'orchestre}
\textbf{l'orchestre} — \textquotedblleft{}orchestra; stalls (in a theatre)\textquotedblright{}.

\begin{quote}
L'orchestre joue ce soir au théâtre.
\end{quote}

\begin{quote}
\emph{The orchestra is playing at the theatre tonight.}
\end{quote}

Masculine. In a theatre, \emph{l'orchestre} is also the stalls.

\begin{cousinweb}
\emph{l'orchestre} comes from Latin \emph{orchēstra}, a borrowing from Greek \emph{orkh\'{\={e}}stra} \textquotedblleft{}the space where the chorus danced\textquotedblright{} (EtymDB), built on \emph{orkheîsthai} \textquotedblleft{}to dance\textquotedblright{}.
\end{cousinweb}

\begin{grammarlens}[title={What you've built}]
One more word for telling, explaining and describing at B1.
\end{grammarlens}

\subsection*{Your turn}
\begin{itemize}
\raggedright
  \item \emph{Say it:} \emph{l'orchestre}
  \item \emph{Say it:} \emph{l'orchestre}, once more
  \item \emph{Say it:} a sentence of your own with \emph{l'orchestre}
  \item \emph{Recall:} say the French for chorus, then the French for melody, then say \emph{l'orchestre} again
\end{itemize}

\begin{tcolorbox}[breakable,colback=teal!4,colframe=teal!35!black,title={Before you move on}]
What does l'orchestre mean? (\textquotedblleft{}Orchestra; stalls (in a theatre)\textquotedblright{}.) Say it once more.
\end{tcolorbox}
\section[le chef d'orchestre]{le chef d'orchestre — conductor (of an orchestra)}

\label{lesson:FR-C513-chef-d-orchestre}



\begin{quote}
Before the new one: say the French for melody, then the French for orchestra.
\end{quote}

\subsection*{You'll want to know: le chef d'orchestre}
\textbf{le chef d'orchestre} — \textquotedblleft{}conductor (of an orchestra)\textquotedblright{}.

\begin{quote}
Le chef d'orchestre lève sa baguette.
\end{quote}

\begin{quote}
\emph{The conductor raises his baton.}
\end{quote}

A woman is often \emph{la cheffe d'orchestre} or \emph{la chef d'orchestre}.

\begin{cousinweb}
\emph{le chef d'orchestre} is \textquotedblleft{}the head of the orchestra\textquotedblright{}: \emph{chef} is inherited from Latin \emph{caput} \textquotedblleft{}head\textquotedblright{} (Lewis \& Short).
\end{cousinweb}

\begin{grammarlens}[title={What you've built}]
One more word for telling, explaining and describing at B1.
\end{grammarlens}

\subsection*{Your turn}
\begin{itemize}
\raggedright
  \item \emph{Say it:} \emph{le chef d'orchestre}
  \item \emph{Say it:} \emph{le chef d'orchestre}, once more
  \item \emph{Say it:} a sentence of your own with \emph{le chef d'orchestre}
  \item \emph{Recall:} say the French for melody, then the French for orchestra, then say \emph{le chef d'orchestre} again
\end{itemize}

\begin{tcolorbox}[breakable,colback=teal!4,colframe=teal!35!black,title={Before you move on}]
What does le chef d'orchestre mean? (\textquotedblleft{}Conductor (of an orchestra)\textquotedblright{}.) Say it once more.
\end{tcolorbox}
\section[la partition]{la partition — (musical) score, sheet music}

\label{lesson:FR-C513-partition}



\begin{quote}
Before the new one: say the French for orchestra, then the French for conductor.
\end{quote}

\subsection*{You'll want to know: la partition}
\textbf{la partition} — \textquotedblleft{}(musical) score, sheet music\textquotedblright{}.

\begin{quote}
Le pianiste a oublié sa partition à la maison.
\end{quote}

\begin{quote}
\emph{The pianist left his sheet music at home.}
\end{quote}

Also a division, as in English \emph{partition}.

\begin{cousinweb}
\emph{la partition} is a learned borrowing from Latin \emph{partītiō} \textquotedblleft{}a division\textquotedblright{} (EtymDB), from \emph{partīrī} \textquotedblleft{}to divide\textquotedblright{}, built on \emph{pars} \textquotedblleft{}part\textquotedblright{} (Lewis \& Short).
\end{cousinweb}

\begin{grammarlens}[title={What you've built}]
One more word for telling, explaining and describing at B1.
\end{grammarlens}

\subsection*{Your turn}
\begin{itemize}
\raggedright
  \item \emph{Say it:} \emph{la partition}
  \item \emph{Say it:} \emph{la partition}, once more
  \item \emph{Say it:} a sentence of your own with \emph{la partition}
  \item \emph{Recall:} say the French for orchestra, then the French for conductor, then say \emph{la partition} again
\end{itemize}

\begin{tcolorbox}[breakable,colback=teal!4,colframe=teal!35!black,title={Before you move on}]
What does la partition mean? (\textquotedblleft{}(musical) score, sheet music\textquotedblright{}.) Say it once more.
\end{tcolorbox}
\section[la répétition]{la répétition — rehearsal; repetition}

\label{lesson:FR-C513-repetition}



\begin{quote}
Before the new one: say the French for conductor, then the French for (musical) score.
\end{quote}

\subsection*{You'll want to know: la répétition}
\textbf{la répétition} — \textquotedblleft{}rehearsal; repetition\textquotedblright{}.

\begin{quote}
La répétition commence à vingt heures.
\end{quote}

\begin{quote}
\emph{The rehearsal starts at eight o'clock.}
\end{quote}

\emph{Une répétition générale} — a dress rehearsal.

\begin{cousinweb}
\emph{la répétition} is a learned borrowing from Latin \emph{repetītiō} \textquotedblleft{}a repeating\textquotedblright{} (EtymDB), from \emph{repetere} \textquotedblleft{}to go back to\textquotedblright{}, built on \emph{petere} \textquotedblleft{}to seek\textquotedblright{} (Lewis \& Short).
\end{cousinweb}

\begin{grammarlens}[title={What you've built}]
One more word for telling, explaining and describing at B1.
\end{grammarlens}

\subsection*{Your turn}
\begin{itemize}
\raggedright
  \item \emph{Say it:} \emph{la répétition}
  \item \emph{Say it:} \emph{la répétition}, once more
  \item \emph{Say it:} a sentence of your own with \emph{la répétition}
  \item \emph{Recall:} say the French for conductor, then the French for (musical) score, then say \emph{la répétition} again
\end{itemize}

\begin{tcolorbox}[breakable,colback=teal!4,colframe=teal!35!black,title={Before you move on}]
What does la répétition mean? (\textquotedblleft{}Rehearsal; repetition\textquotedblright{}.) Say it once more.
\end{tcolorbox}
\section[improviser]{improviser — to improvise}

\label{lesson:FR-C513-improviser}



\begin{quote}
Before the new one: say the French for (musical) score, then the French for rehearsal.
\end{quote}

\subsection*{You'll want to know: improviser}
\textbf{improviser} — \textquotedblleft{}to improvise\textquotedblright{}.

\begin{quote}
Les musiciens de jazz improvisent beaucoup.
\end{quote}

\begin{quote}
\emph{Jazz musicians improvise a lot.}
\end{quote}

Also outside music: \emph{on a improvisé un repas}.

\begin{cousinweb}
\emph{improviser} was borrowed from Italian \emph{improvvisare} (EtymDB), from \emph{improvviso}, from Latin \emph{imprōvīsus} \textquotedblleft{}unforeseen\textquotedblright{}, built on \emph{vidēre} \textquotedblleft{}to see\textquotedblright{} (Lewis \& Short).
\end{cousinweb}

\begin{grammarlens}[title={What you've built}]
One more word for telling, explaining and describing at B1.
\end{grammarlens}

\subsection*{Your turn}
\begin{itemize}
\raggedright
  \item \emph{Say it:} \emph{improviser}
  \item \emph{Say it:} \emph{improviser}, once more
  \item \emph{Say it:} a sentence of your own with \emph{improviser}
  \item \emph{Recall:} say the French for (musical) score, then the French for rehearsal, then say \emph{improviser} again
\end{itemize}

\begin{tcolorbox}[breakable,colback=teal!4,colframe=teal!35!black,title={Before you move on}]
What does improviser mean? (\textquotedblleft{}To improvise\textquotedblright{}.) Say it once more.
\end{tcolorbox}
\section[(dialogue)]{Retour sur les chapitres 506–513}

\label{lesson:FR-R513-revision-506-513}



\begin{quote}
Say the French for draw and ranking.
\end{quote}

\subsection*{Your turn}
Say these aloud:

\begin{itemize}
\raggedright
  \item draw, ranking, to knock out, (se qualifier) to qualify, winner
  \item tournament, medal, podium, champion, to face
  \item runner, finishing line, marathon, relay, to time
  \item to dive, dive, swimsuit, breaststroke, to row
  \item skate, to skate, sledge, ice rink, slope
  \item flute, cello, trumpet, accordion, to tune
  \item choir, to hum, to whistle softly, chorus, melody
  \item orchestra, conductor, (musical) score, rehearsal, to improvise
\end{itemize}

\begin{tcolorbox}[breakable,colback=teal!4,colframe=teal!35!black,title={Before you move on}]
Draw? (\textbf{le match nul}.) Ranking? (\textbf{le classement}.) To knock out? (\textbf{éliminer}.) (se qualifier) to qualify? (\textbf{qualifier}.) Winner? (\textbf{le vainqueur}.) Tournament? (\textbf{le tournoi}.) Medal? (\textbf{la médaille}.) Podium? (\textbf{le podium}.) Champion? (\textbf{le champion}.) To face? (\textbf{affronter}.) Runner? (\textbf{le coureur}.) Finishing line? (\textbf{la ligne d'arrivée}.) Marathon? (\textbf{le marathon}.) Relay? (\textbf{le relais}.) To time? (\textbf{chronométrer}.) To dive? (\textbf{plonger}.) Dive? (\textbf{le plongeon}.) Swimsuit? (\textbf{le maillot de bain}.) Breaststroke? (\textbf{la brasse}.) To row? (\textbf{ramer}.) Skate? (\textbf{le patin}.) To skate? (\textbf{patiner}.) Sledge? (\textbf{la luge}.) Ice rink? (\textbf{la patinoire}.) Slope? (\textbf{la piste}.) Flute? (\textbf{la flûte}.) Cello? (\textbf{le violoncelle}.) Trumpet? (\textbf{la trompette}.) Accordion? (\textbf{l'accordéon}.) To tune? (\textbf{accorder}.) Choir? (\textbf{la chorale}.) To hum? (\textbf{fredonner}.) To whistle softly? (\textbf{siffloter}.) Chorus? (\textbf{le refrain}.) Melody? (\textbf{la mélodie}.) Orchestra? (\textbf{l'orchestre}.) Conductor? (\textbf{le chef d'orchestre}.) (musical) score? (\textbf{la partition}.) Rehearsal? (\textbf{la répétition}.) To improvise? (\textbf{improviser}.)
\end{tcolorbox}
